\section{System Design}

\begin{figure}[t]
\centering
\begin{tikzpicture}[
  font=\small,
  node distance=6mm,
  box/.style={draw=oaink, rounded corners=2pt, align=center, inner sep=4pt,
    minimum height=8mm, fill=white},
  stage/.style={box, fill=oagray!50, minimum width=20mm},
  gate/.style={draw=oablue, rounded corners=2pt, align=center, inner sep=3pt,
    fill=oablue!8, font=\footnotesize},
  outcome/.style={align=center, font=\footnotesize\bfseries, inner sep=2pt},
  flow/.style={-{Stealth[length=2mm]}, draw=oaink, thick},
  ]
  \node[stage] (rec) {\textbf{Record}\\one demo};
  \node[stage, right=of rec] (comp) {\textbf{Compile}\\bundle};
  \node[stage, right=of comp] (rep) {\textbf{Replay}\\deterministic};
  \draw[flow] (rec) -- (comp);
  \draw[flow] (comp) -- (rep);

  % Governance gate stack, to the right of replay (left-aligned so the widest
  % gate does not overrun the Replay node).
  \coordinate (gx) at ($(rep.east)+(15mm,0)$);
  \node[gate, anchor=west] (id)   at ($(gx)+(0,9mm)$)   {identity};
  \node[gate, anchor=west] (post) at ($(gx)+(0,3mm)$)   {postcondition};
  \node[gate, anchor=west] (eff)  at ($(gx)+(0,-3mm)$)  {effect (system of record)};
  \node[gate, anchor=west] (pol)  at ($(gx)+(0,-9mm)$)  {policy};
  \begin{scope}[on background layer]
    \node[draw=oablue, dashed, rounded corners=3pt, fit=(id)(post)(eff)(pol),
      inner sep=3pt, label={[font=\footnotesize\itshape,oablue]above:governance gates}] (gatebox) {};
  \end{scope}
  \draw[flow] (rep) -- (gatebox.west);

  % Three outcomes.
  \node[outcome, right=10mm of gatebox, yshift=8mm, text=oagreen] (ok) {Success\\(verified)};
  \node[outcome, below=3mm of ok, text=oaamber] (heal) {Repair\\(logged)};
  \node[outcome, below=3mm of heal, text=oared] (halt) {Refuse\\(halt)};
  \draw[flow, oagreen] (gatebox.east) -- (ok);
  \draw[flow, oaamber] (gatebox.east) -- (heal);
  \draw[flow, oared] (gatebox.east) -- (halt);

  % Repair feeds back a patch; refusal stops.
  \draw[flow, oaamber, dashed] (heal.south) to[out=-90,in=-90] node[below,font=\scriptsize\itshape]{reviewable patch} (rep.south);
\end{tikzpicture}
\caption{\system{} compiles one demonstration into a deterministic program.
Healthy replay makes no model calls. Each step passes governance gates; the
runtime verifies and succeeds, records a bounded \emph{repair} when a lower-
fidelity rung re-resolves a changed target, or \emph{refuses} when a configured
check cannot be verified. A repair is an auditable patch, never a silent
overwrite of the deployed bundle.}
\label{fig:pipeline}
\end{figure}

\subsection{Record and compile}
A recorder captures actions, frames, element observations, and structural state
when available. Declared secret fields are replaced with runtime parameters.
The compiler emits a bundle containing a workflow program, target crops,
locators, landmarks, postconditions, parameter declarations, and a manifest.
The linear trace is the simplest program; the intermediate representation also
supports states, guarded transitions, branches, loops, and subflows. Given
several demonstrations of the same task, a deterministic multi-trace induction
pass promotes a value that varies at an aligned position to a typed parameter, a
consecutive repeated sub-sequence whose repetition count varies to a loop over an
inferred relation, and a divergence under a detectable condition to a guarded
branch; it makes no model call. Its reach is deliberately narrow and we state
the limit rather than the aspiration: only consecutive repeated bodies are
recognized, so a navigate-away-and-return loop, a paginated loop, and a
per-row-conditional body are not. Where induction cannot certify the shape---an
ambiguous loop around an irreversible action, an undetermined branch condition,
an unhandled zero-result case---it emits an uncertainty and quarantines the
result rather than inventing one, and the compiled workflow stays uncertified
until an operator answers a concrete, grounded question about the specific
divergence.

Parameterization is where a compiler is most tempted to guess, because the
difference between a constant and a parameter is the difference between
replaying a demonstration and running a new case. \system{} therefore infers
parameter \emph{candidates} deterministically---a value typed into a field
carries the field's own label, taken from the recorder's structural capture or,
failing that, from the nearest OCR line left of or above the field rectangle---
and then refuses to apply them. Each candidate is written to a sidecar with the
example value masked, and one batch operator pass confirms, renames, marks
secret, or keeps each as a constant. The default on an unanswered prompt is
\emph{keep constant}: an unconfirmed inference changes nothing, and the bundle
replays exactly as demonstrated. This is deliberate rather than conservative
for its own sake---the parameter name is the identifier an effect contract binds
to and that secrets are keyed by, so naming is confirmed by a human, never
inferred behind the operator's back. The inference makes no model call.

\subsection{Resolution ladder}
For an action target, replay tries evidence in descending order of fidelity:
structured DOM or accessibility identity, local and global template matching,
OCR, landmark geometry, and an optional grounding model. A backend presents a
small common contract for observation and actuation. This provides a pixel
floor without discarding structured evidence on browsers and native desktops.

Healthy execution stops at deterministic tiers and makes zero model calls. A
lower rung resolving a changed but verified target creates a repair patch. The
runtime can save a healed derivative bundle, but does not silently replace the
deployed source bundle.

\begin{figure}[t]
\centering
\begin{tikzpicture}[
  font=\footnotesize,
  rung/.style={draw=oaink, rounded corners=2pt, minimum width=62mm,
    minimum height=6.5mm, align=center, inner sep=2pt},
  ]
  \node[rung, fill=oagreen!14] (r1) {1. Structured DOM / accessibility identity};
  \node[rung, fill=oagreen!10, below=1.6mm of r1] (r2) {2. Local template match (cropped target)};
  \node[rung, fill=oagreen!7, below=1.6mm of r2] (r3) {3. Global template match};
  \node[rung, fill=oaamber!12, below=1.6mm of r3] (r4) {4. OCR text anchor};
  \node[rung, fill=oaamber!16, below=1.6mm of r4] (r5) {5. Landmark geometry \quad (pixel floor)};
  \node[rung, fill=oared!12, below=1.6mm of r5] (r6) {6. Grounding model \quad (optional, off by default)};

  \draw[-{Stealth[length=2mm]}, very thick, oaink]
    ($(r1.north west)+(-6mm,0)$) -- ($(r6.south west)+(-6mm,0)$)
    node[midway, rotate=90, anchor=south, font=\footnotesize\itshape]{descending fidelity};

  \node[right=3mm of r3, align=left, font=\scriptsize, text=oagreen!60!black] (det)
    {deterministic;\\zero model calls};
  \draw[decorate, decoration={brace, amplitude=4pt}, oagreen!60!black]
    ($(r1.north east)+(1mm,0)$) -- ($(r5.south east)+(1mm,0)$);
  \node[right=6mm of r4, font=\scriptsize\itshape, text=oared] {model calls only here};
\end{tikzpicture}
\caption{The resolution ladder tries evidence in descending order of fidelity.
Structured identity is preferred; local visual evidence gives a pixel floor for
substrates without a DOM. Rungs 1--5 are deterministic and make no model calls.
A grounding model is an optional last rung, disabled by default, and its rescue
is subject to the same identity and effect checks.}
\label{fig:ladder}
\end{figure}

\subsection{Postconditions and reports}
Each step may carry visual or structural postconditions such as text presence,
region stability, URL change, or window state. Replay writes a human-readable
and machine-readable report with rung, confidence, identity verdict,
postconditions, effects, model calls, elapsed time, repair, and halt reason.
Postconditions are necessary but not sufficient: a success banner is evidence
about the screen, not independent evidence about a database write.

\subsection{Backends and explicit surface selection}
The browser backend is the reference implementation and supplies published
end-to-end evidence against a real third-party application. Windows UI
Automation adds structured candidate enumeration and native actions behind a
typed, authenticated in-session contract. Native macOS binds capture and input
to a unique application window. The RDP and Citrix Workspace backends expose the
same runtime contract over a pixel-only network session. Section~4 reports
scoped task and effect qualifications for the Windows, macOS, and RDP
mechanisms, and a deterministic synthetic stand-in campaign for the Citrix
Workspace backend. A real Citrix ICA/HDX environment is a separate deployment
and inherits neither the RDP qualification nor the stand-in campaign.

The browser is one execution surface among six---\texttt{web},
\texttt{windows}, \texttt{macos}, \texttt{linux}, \texttt{rdp}, and
\texttt{citrix}---and not a privileged default. Each surface fixes an execution
\emph{mode}: \texttt{rdp} and \texttt{citrix} are \emph{external}, driving a
local client window of a remote session through pixels, keyboard, and mouse with
nothing installed inside that session; the others are \emph{in-session}. The
mode is fixed by capability negotiation at qualification and is never silently
switched at run time. Two consequences are enforced at the command line rather
than documented as advice. First, under the production execution profiles a
record or run command that does not name a surface is \emph{refused}, because an
implicit browser default silently decides the substrate a workflow will be
qualified on. Second, a compiled bundle carries the surface it was recorded on,
and running it on a different surface is refused unless the operator passes an
explicit override, which is then recorded in the run report as compatibility
evidence. A cross-surface run is never silent.
